\documentclass[11pt]{article}
\usepackage[T1]{fontenc}
\usepackage{amsmath,amssymb,amsthm}
\usepackage{geometry}
\geometry{margin=1in}
\newtheorem{theorem}{Theorem}
\newtheorem{proposition}{Proposition}
\newcommand{\Ftwo}{\mathbb F_2}
\newcommand{\Gproj}{G_{\mathrm{proj}}}
\title{Laboratory V73: Bicriteria Residual Ordering and Counted Branch Decompositions}
\author{Internal research artifact; authorship metadata pending}
\date{31 July 2026}
\begin{document}
\maketitle

\begin{abstract}
We give an exact subset dynamic program minimizing projected affine-residual
cost under a support-frontier budget, augment branch residual states with exact
multiplicities, identify the precise supplied-target avoidance certificate
provided by root multiplicity zero, and prove that the existing normalized
positive-fiber schema cannot itself produce an avoided target.  We also prove
that the private-vertex partition-zero binary-tree family has unbounded support
linear width but optimum projected residual cost exactly equal to the number of
gates.  No unrestricted Range Avoidance algorithm or P-versus-NP consequence
is claimed.
\end{abstract}

\section{Budgeted residual ordering}

For a processed gate set $S$, let $w(S)$ be the number of nonempty projected
affine residuals and let $\lambda(S)$ be the support frontier.  Under the
repository convention,
\[
  \Gproj(\pi)=\sum_{i=0}^{m-1} w(P_i),
\]
where $P_0=\varnothing$ and $P_i$ is the first $i$ gates of $\pi$.

\begin{theorem}[Exact bicriteria dynamic program]
Fix a frontier budget $B$.  The minimum projected residual cost among orders
whose every prefix has frontier at most $B$ is given by
\[
 C_B(\varnothing)=0,\qquad
 C_B(S)=\min_{e\in S}\left(C_B(S\setminus\{e\})+w(S\setminus\{e\})\right),
\]
with transitions restricted to feasible predecessors and $\lambda(S)\le B$.
Given the subset tables, the optimization uses $O(m2^m)$ transitions and
$O(2^m)$ stored costs.
\end{theorem}

\begin{proof}
In any optimum order of $S$, remove the last gate $e$.  The preceding order is
feasible for $S\setminus\{e\}$ and the final transition contributes
$w(S\setminus\{e\})$, proving the lower bound.  Conversely, append $e$ to an
optimum feasible predecessor.  The frontier test on $S$ proves feasibility and
the matching upper bound.
\end{proof}

\section{Multiplicity-enhanced branch residuals}

For a node $t$ of a binary gate decomposition, let $\partial t$ be the
variables shared by its leaf set and the complementary gates.  For every
canonical residual $A$ on $\partial t$, let $\mu_t(A)$ count complete cell
selections in the subtree that are consistent and project to $A$.

\begin{theorem}[Exact branch multiplicities]
At a join with children $u,v$, if the consistent intersection of residuals
$A,B$ projects to $R$ on the parent boundary, then adding
\[
  \mu_u(A)\mu_v(B)
\]
to $\mu_t(R)$ computes all multiplicities exactly.
\end{theorem}

\begin{proof}
Every complete cell selection splits uniquely into one selection in each
child.  Canonical projection assigns each child selection to exactly one
residual.  Products count Cartesian pairs and sums combine disjoint selections
that project to the same parent residual.  Induction on the tree completes the
proof.
\end{proof}

\begin{proposition}[Supplied-target certificate]
Suppose an output word $y$ is supplied and every fiber $f_i^{-1}(y_i)$ is
represented as an exact disjoint union of affine cells.  Then the root
multiplicity is zero if and only if $y$ is outside the circuit image.
\end{proposition}

\begin{proof}
A complete cell selection is consistent exactly when some input belongs to
every selected gate fiber, equivalently when some input maps to $y$.
\end{proof}

\begin{proposition}[Normalized-schema barrier]
In the current partition-$0,1,2$ compiler, cell zero of every gate contains the
all-zero assignment.  Consequently every normalized instance has a consistent
all-zero complete branch, so its root multiplicity is at least one.
\end{proposition}

Thus the current schema certifies a supplied target only after both output
polarities and exact fiber representations are provided; it cannot search for
an avoided target by itself.

\section{Binary-tree compression}

Orient each tree edge from parent $p$ to child $c$, add a private vertex $z$,
and use cells
\[
  p=c=z=0,\qquad p=0,\ c+z=1.
\]

\begin{theorem}[Unbounded width with minimum residual cost]
For the private-vertex partition-zero tree family, rooted postorder has one
projected residual at every nonterminal layer.  Therefore
\[
  \Gproj^*=m.
\]
For a tree-aligned branch decomposition, every gate leaf has at most two
residual states and every internal node has exactly one.
\end{theorem}

\begin{proof}
A completed child subtree projects to the single condition that its top
vertex is zero.  Joining its incoming edge eliminates the child and private
vertex and leaves only the parent-zero condition.  Induction along postorder
gives one residual per layer, hence an upper bound $m$.  Every one of the $m$
nonterminal layers is nonempty, giving the matching lower bound.  The same
induction gives the branch-state statement; a bare edge to an unfinished
internal child is the only place where two residuals can occur.
\end{proof}

V72 proved that the same support family has primal treewidth at most two and
unbounded support linear width.  Hence support-width growth does not imply
growth of $\Gproj^*$ on this family.

\section{Nonclaims}
The results are internally verified but not peer reviewed or novelty
confirmed.  They do not prove unrestricted $\mathrm{NC}^0_3$ Range Avoidance,
an all-orders superpolynomial lower bound, a simulation to a standard proof or
branching-program model, or any consequence for P versus NP.

\end{document}
